\chapter{Validation, Defect Calculus, and Method Promotion}
\label{ch:bf-highdim-validation-defect-promotion}
\label{ch:bf-highdim-candidate-synthesis}

This chapter closes the restarted high-dimensional block.  It synthesizes the
validation and promotion logic shared by the deterministic Gaussian/sparse-grid
lane and the non-Gaussian low-rank retained-object lane.  Its purpose is not to
re-derive the methods, but to consume their exported defects, diagnostics,
approximate-target contracts, and HMC-admissibility consequences.

The previous chapters have already separated:
\begin{itemize}[leftmargin=0.28in]
  \item exact target from approximate target,
  \item carried object from moment or retained-object engine,
  \item frozen branch-defining structure from recomputed numerical values.
\end{itemize}
The role here is therefore downstream and comparative: which diagnostics veto a
method, which diagnostics only explain it, and under what conditions a chapter-
local approximation becomes an industrially credible component rather than a
merely interesting computation.

\section{Validation Architecture And Benchmark Roles}
\label{sec:bf-hd-validation-architecture}

The restarted high-dimensional block uses three ledgers of evidence that must not
be conflated:
\begin{enumerate}[label=(\arabic*)]
  \item \textbf{Engineering correctness.}  Did the declared scalar build, factor,
  and update without violating its own branch rules?
  \item \textbf{Numerical or sampler validity.}  Did finite-difference parity,
  ESS, PSD, rank, support, Jacobian, or divergence diagnostics pass?
  \item \textbf{Scientific interpretation.}  Given the surviving diagnostics,
  what may be concluded about approximation quality, target fidelity, or method
  promotion?
\end{enumerate}
A result in one ledger cannot be silently upgraded into the next.  A green build
is engineering evidence.  A valid same-scalar finite-difference row is numerical
evidence.  Neither by itself is a scientific proof that the approximation is the
right industrial method.

The benchmark roles are therefore tiered:
\begin{itemize}[leftmargin=0.28in]
  \item exact or controlled references test whether the declared scalar behaves
  as intended on small models,
  \item local diagnostics test whether the approximation object stays meaningful,
  \item promotion ladders decide whether a method deserves more expensive
  downstream comparison.
\end{itemize}

\section{Finite-Difference Verification As The First Shared Gate}
\label{sec:bf-hd-validation-fd-first}

The first shared validation gate is finite-difference parity on a fixed branch.
For a base parameter \(\theta_0\) and coordinate direction \(e_i\), define
\begin{equation}
  D_i(h)=
  \frac{\widehat\ell_T(\theta_0+h e_i;B)-\widehat\ell_T(\theta_0-h e_i;B)}{2h},
  \qquad
  G_i=\partial_i\widehat\ell_T(\theta_0;B),
  \label{eq:bf-hd-validation-fd}
\end{equation}
where both perturbed evaluations reuse the same fixed branch ledger and
recompute the parameter-dependent numerical objects through the same declared
equations.  The row is branch-valid only if the branch identity is unchanged on
both sides.

A useful reporting contract is:
\begin{equation}
  e_i(h)=\frac{|D_i(h)-G_i|}{\max\{1,|D_i(h)|,|G_i|\}},
  \label{eq:bf-hd-validation-fd-error}
\end{equation}
with a small ladder such as
\begin{equation}
  h\in\{10^{-1},10^{-2},10^{-3},10^{-4},10^{-5}\}\max\{1,|\theta_i|\}.
  \label{eq:bf-hd-validation-fd-ladder}
\end{equation}
The diagnostic supports the same-scalar claim only when branch-valid rows show a
stable decreasing error window.  Failure may indicate a wrong derivative, a
branch mismatch, or a nonsmooth/ill-conditioned value path.  It is a validation
gate, not a posterior-accuracy theorem.

\section{Why The Synthesis Must Start With Defects}
\label{sec:bf-hd-synthesis-purpose}

The preceding chapters explain high-dimensional nonlinear filtering methods.
An industrial macro-finance program needs a harsher object: a calculus of
failure.  A method is useful for a global-bank or central-bank model only when
the failure mode is explicit, the mitigation has a mathematical contract, and
the cost variable that can kill the method is visible before implementation.
Academic novelty is not the target here.  The target is a defensible
composition of known methods under stated assumptions.

This chapter therefore treats every candidate as a conditional component.  A
Gaussian scaffold is useful if it exposes local curvature and covariance
failure.  Sparse-grid or high-degree cubature is useful if it diagnoses local
nonlinearity without pretending to be a global filter.  Tensor methods are
useful if economic structure keeps ranks stable and preserves probability or
covariance semantics.  Transport maps are useful if they reduce geometry while
leaving a checkable target.  HMC is useful only downstream of a same-scalar
likelihood contract.

The reader-facing synthesis map is this: start with the exact conditional law
and likelihood normalizers from Chapter~\ref{ch:bf-highdim-nonlinear-foundations};
replace them only by named approximations; attach one diagnostic to each object
being approximated; and promote a method only after the previous object exports
the input required by the next.  The chapter's exact target is not a new
algorithm.  It is an auditable decision procedure for composing existing methods
without losing the scalar likelihood, probability-law, support, Jacobian, rank,
or covariance semantics needed by industrial inference.

\paragraph{Running quadratic-observation cell.}
The scalar cell
\[
  x_t=\rho x_{t-1}+\eta_t,\qquad y_t=x_t^2+\epsilon_t
\]
has now served five roles.  Chapter~\ref{ch:bf-highdim-nonlinear-foundations}
used it to show that the exact filtering law can have two lobes.
Chapter~\ref{ch:bf-highdim-gaussian-quadrature} used it to show why Gaussian
projection can preserve moments while losing shape.  Chapter~\ref{ch:bf-highdim-fixed-cloud-sgqf-validation}
used it to turn one fixed sparse-grid cloud into a concrete value-path and
numeric-oracle contract.  Chapter~\ref{ch:bf-highdim-low-rank-density-kr} and
Chapter~\ref{ch:bf-highdim-squared-tt-fixed-branch-likelihoods} used it to
separate retained-object semantics from fixed-branch likelihood construction.
Chapter~\ref{ch:bf-highdim-hmc-research} used it to insist that HMC needs a
scalar likelihood and a same-scalar gradient.  Here the same cell is not
evidence for a macro-finance model.  It is the small mechanism that tells a
mixed numerical reader what the architecture is trying not to lose.

\paragraph{Imports from the preceding chapters.}
\begingroup
\scriptsize
\setlength{\tabcolsep}{2pt}
\begin{longtable}{@{}p{0.14\linewidth}p{0.19\linewidth}p{0.22\linewidth}p{0.17\linewidth}p{0.18\linewidth}@{}}
\toprule
Source chapter & Exported object & Exported defect & Exported diagnostic/veto
& Consumed here as \\
\midrule
\endhead
Ch.~\ref{ch:bf-highdim-nonlinear-foundations} & conditional law, normalizer,
likelihood gradient & wrong target, missing normalizer, value-gradient
mismatch & finite \(Z_t\), mass, positivity, same-scalar check & target and
likelihood gates \\
Ch.~\ref{ch:bf-highdim-gaussian-quadrature} & Gaussian moment pair and
point-rule family & projection bias, exactness without accuracy, point explosion
& innovation, PSD, active dimension, comparator residual & scaffold and local
curvature diagnostic \\
Ch.~\ref{ch:bf-highdim-fixed-cloud-sgqf-validation} & fixed sparse-grid scalar,
numeric oracle, branch record & merge mismatch, signed-weight instability,
branch-invalid finite differences & oracle reproduction, branch-valid FD,
minimum-eigenvalue veto, cloud reproduction & deterministic Gaussian lane
validation \\
Ch.~\ref{ch:bf-highdim-low-rank-density-kr} & retained object, coordinate map,
KR / preconditioning semantics & support failure, Jacobian mismatch,
rank/semantic failure & support/logdet, rank, mass, positivity & retained-object
meaning gate \\
Ch.~\ref{ch:bf-highdim-squared-tt-fixed-branch-likelihoods} & fixed-branch TT/KR
scalar, stored-field contract & wrong saved object, branch-identity loss,
incomplete next-step closure & stored-field audit, retained-marginal closure,
branch-fixed structural check & operational middle gate \\
Ch.~\ref{ch:bf-highdim-hmc-research} & scalar potential, gradient, transformed
target & same-scalar mismatch, divergences, invalid transform & scalar parity,
finite support/Jacobian, energy diagnostics & downstream sampler gate \\
\bottomrule
\end{longtable}
\endgroup

\section{Lane-Specific Validation Ladders Before Promotion}
\label{sec:bf-hd-validation-lane-ladders}

The high-dimensional part now contains two distinct approximation lanes with
different local failure modes, so the validation ladder must stay lane-aware
before the closing promotion comparison begins.

\subsection{FixedSGQF ladder}
\label{sec:bf-hd-validation-sgqf-ladder}

The sparse-grid lane should not be promoted from elegant construction alone.
Before any cross-method comparison, it should pass a ladder that includes:
\begin{enumerate}[leftmargin=0.28in]
  \item deterministic cloud reproduction and duplicate-merge checks,
  \item the one-step numeric oracle from
  Chapter~\ref{ch:bf-highdim-fixed-cloud-sgqf-validation},
  \item branch-valid finite differences on a saved fixed-cloud branch,
  \item signed-weight stability checks through \(P_t^-\), \(S_t\), and carried
  covariance objects,
  \item controlled GHQ / linear-Gaussian comparators where available.
\end{enumerate}
These diagnostics answer whether the implementation is evaluating the declared
fixed-cloud scalar, not whether Gaussian projection has become an exact
high-dimensional posterior method.

\subsection{Retained-object TT/KR ladder}
\label{sec:bf-hd-validation-ttkr-ladder}

The TT/KR lane carries a richer retained object and therefore faces richer local
failure modes.  Before cross-method promotion, it should pass a ladder that
includes:
\begin{enumerate}[leftmargin=0.28in]
  \item coordinate-system and Jacobian-consistency checks,
  \item positivity / nonnegativity and normalization checks for the represented
  object,
  \item rank-growth and conditioning diagnostics,
  \item stored-field and retained-marginal closure checks for the next step,
  \item branch-valid finite differences on the fixed-branch likelihood route.
\end{enumerate}
These diagnostics do not prove that moderate TT rank is always sufficient or that
the adaptive source algorithm is globally smooth.  They only establish whether
the frozen retained-object branch remains a scientifically interpretable scalar
object on the tested route.

\paragraph{Copied-core and branch-status reporting.}
When the derivative claim is part of the result, the TT/KR lane should report not
only a finite-difference error table but also the branch-status fields imported
from Chapter~\ref{ch:bf-highdim-hmc-research}: branch-identity indicators,
copied-core indicators, decreasing-window indicators, and the resulting status
classification.  A small error alone is insufficient if the perturbed values were
computed on different frozen fields or with copied fitted cores.

A useful derivative reporting object is:
\begin{equation}
  \mathcal R_{\rm FD}=
  \{\widehat\ell_T(\theta_0;B),\ G_i,\ (h_i,D_i(h_i),e_i(h_i))_i,
    I_\pm(h_i),K_\pm(h_i),W_i,\text{status}\}.
  \label{eq:bf-hd-validation-fd-report-object}
\end{equation}
Here the report should also retain the scalar values
\(\widehat\ell_T(\theta_0+h_i e_j;B)\) and
\(\widehat\ell_T(\theta_0-h_i e_j;B)\) for each \(h_i\), not merely the final
error summary.  The derivative table is meaningful only when the branch record is
identical in the three columns base / plus / minus.

\paragraph{Benchmark-specific interpretation.}
The benchmark ledger should stay explicit about what each model can and cannot
establish.  A linear-Gaussian benchmark can establish exact-reference agreement
for the declared run contract, but not nonlinear posterior accuracy.  A
stochastic-volatility or predator-prey benchmark can establish long-horizon
nonlinear stability or preconditioning benefit under the declared run contract,
but not universal scalability.  A spatial SIR benchmark can establish whether the
method remains useful at the tested state dimension, horizon, rank ladder, and
basis ladder, but not that all spatial epidemic models admit moderate TT rank.
One-seed demonstrations remain diagnostic unless the result artifact explicitly
justifies a stronger interpretation.

\paragraph{Runtime, memory, and veto interpretation.}
Runtime and memory should be interpreted only after veto diagnostics have passed.
When reported, they should be split into mathematically meaningful components:
forward target evaluation, fitting solves, mass contractions, KR inversion or
sampling work, retained-object storage work, and derivative overhead.  Likewise,
peak memory should separate value-only storage from value-plus-gradient storage.
These are not mere profiling counters; they are the places where the fixed-branch
approximation can become unaffordable or mathematically fragile.

\section{Mathematical Veto Ledger Before Interpretation}
\label{sec:bf-hd-validation-veto-ledger}

The validation report should mark a run as failed before interpreting posterior
plots, trajectory summaries, or method rankings if any mathematical veto fires.
Let
\begin{equation}
  \mathcal F_t=
  \{Z_t,\widehat Z_t,\widehat p_t,C_{t,k},\dot C_{t,k},w_i,D_h,G_{\rm FB}\}
  \label{eq:bf-hd-veto-ledger-objects}
\end{equation}
collect the scalar, tensor, weight, and derivative objects used at time \(t\).
The finite-value veto is
\begin{equation}
  V_{\rm finite}=
  \mathbf{1}\{\exists t,\exists a\in\mathcal F_t:\ a\text{ contains NaN or Inf}\}.
  \label{eq:bf-hd-veto-finite}
\end{equation}
The normalization veto is
\begin{equation}
  V_{\rm norm}=
  \mathbf{1}\{\exists t:\ \widehat Z_t\le0\ \text{or}\ |\int \widehat p_t(z)\dd z-1|>\varepsilon_{\rm norm}\}.
  \label{eq:bf-hd-veto-norm}
\end{equation}
The least-squares conditioning veto is
\begin{equation}
  V_{\rm cond}=
  \mathbf{1}\{\max_{t,k}\kappa(N_{t,k})>\kappa_{\rm max}^{\rm allowed}\}.
  \label{eq:bf-hd-veto-cond}
\end{equation}
The rank-saturation veto is
\begin{equation}
  V_{\rm rank}=
  \mathbf{1}\{\exists t:\ R_t=R_{\rm max}\ \text{and}\ {\rm residual}_t>\varepsilon_{\rm fit}\}.
  \label{eq:bf-hd-veto-rank}
\end{equation}
The KR monotonicity veto is
\begin{equation}
  V_{\rm KR}=
  \mathbf{1}\{\exists j,z_{<j},a<b:\ F_j(b\mid z_{<j})<F_j(a\mid z_{<j})\}.
  \label{eq:bf-hd-veto-kr}
\end{equation}
Finally, the fixed-branch derivative veto is
\begin{equation}
  V_{\rm grad}=
  \mathbf{1}\{B(\beta+h)\ne B(\beta)\ \text{or}\ B(\beta-h)\ne B(\beta)\ \text{or no branch-stable decreasing window holds}\}.
  \label{eq:bf-hd-veto-grad}
\end{equation}
A run should be interpreted only if
\begin{equation}
  V_{\rm finite}+V_{\rm norm}+V_{\rm cond}+V_{\rm rank}+V_{\rm KR}+V_{\rm grad}^{\rm required}=0,
  \label{eq:bf-hd-veto-sum-rule}
\end{equation}
where \(V_{\rm grad}^{\rm required}=V_{\rm grad}\) only when the claim includes
validating the derivative.  The tolerances in these vetoes must be declared
before a run.  A minimal default table is:
\begin{equation}
\begin{array}{c|c|c}
\text{quantity} & \text{default} & \text{meaning}\\
\hline
\varepsilon_{\rm norm} & 10^2\epsilon_{\rm mach}+10\,\varepsilon_{\rm quad} & \text{normalization tolerance}\\
\kappa_{\rm max}^{\rm allowed} & 10^{10}\ \text{in double precision} & \text{least-squares conditioning ceiling}\\
\varepsilon_{\rm fit} & \max(10^{-8},10\,{\rm tol}_{\rm ALS}) & \text{rank-saturation residual ceiling}\\
\varepsilon_{\rm KR} & 10^2\epsilon_{\rm mach}+10\,\varepsilon_{\rm quad} & \text{monotonicity/inversion tolerance}
\end{array}
  \label{eq:bf-hd-veto-defaults}
\end{equation}
These are implementation tolerances, not scientific outcomes.  Changing them
after seeing posterior plots invalidates the benchmark interpretation.

\section[Particle Collapse]{Defect 1: Particle Collapse Is Exponential Without Structure}
\label{sec:bf-hd-defect-particles}

The bootstrap and SIR filters in Chapter~\ref{ch:bf-highdim-particle-transport-tensor}
are honest baselines because their collapse mechanism is visible.  Let
$W_i=\exp(L_i)$ be unnormalized weights and
$\bar W_i=W_i/\sum_{j=1}^N W_j$.  The effective sample size is
\begin{equation}
  \operatorname{ESS}=\frac{1}{\sum_{i=1}^N \bar W_i^2}.
  \label{eq:bf-hd-synth-ess}
\end{equation}
Suppose the log-weight decomposes across conditionally independent observation
blocks,
\begin{equation}
  L_i=\sum_{k=1}^{d_{\mathrm{eff}}} \ell_{ik},\qquad
  \E[\ell_{ik}]=\mu,\qquad \Var(\ell_{ik})=\sigma_\ell^2,
  \label{eq:bf-hd-synth-logweight}
\end{equation}
with weak dependence sufficient for
\(\Var(L_i)\asymp d_{\mathrm{eff}}\sigma_\ell^2\).  This is the informal
calculus behind the high-dimensional particle-filter warnings of
\citet{bengtsson2008curse} and \citet{snyder2008obstacles}.

\begin{proposition}[Log-weight variance warning]
\label{prop:bf-hd-particle-collapse-calculus}
If the log-weight variance grows like
\(\Var(L_i)=v_d\asymp d_{\mathrm{eff}}\sigma_\ell^2\), then a particle filter
with proposal independent of the current observation can enter a regime where
normalized weights are dominated by a small number of particles unless ensemble
size grows very aggressively with that dispersion scale.
\end{proposition}

\begin{proof}[Derivation]
For a lognormal heuristic, \(\E[W]=\exp(\mu_d+v_d/2)\) and
\(\E[W^2]=\exp(2\mu_d+2v_d)\), so
\[
  \frac{\E[W^2]}{\E[W]^2}=\exp(v_d).
\]
The usual importance-sampling variance factor is therefore exponential in
\(v_d\).  Since the denominator of \eqref{eq:bf-hd-synth-ess} is driven by the
largest normalized weights, stable ESS typically requires sample growth that
offsets this exponential dispersion.  The cited papers give sharper asymptotic
statements under stronger assumptions; the monograph uses this calculation to
expose the industrial defect: without locality, guided proposals, or likelihood
tempering, dimension enters through log-likelihood variance, not through a
benign constant.
\end{proof}

\section{Validation Benchmarks And Promotion Logic}
\label{sec:bf-hd-validation-benchmarks}

The validation ladder should combine exact or controlled references, local
object-validity diagnostics, and promotion rules.  Exact or reference-density
metrics such as Hellinger distance, relative \(L^1\) error, moment error,
trajectory RMSE, coverage, and ESS quantiles are meaningful only when the target
and branch contracts have already been named.  Their role is not to replace the
same-scalar contract, but to test whether a declared approximation remains
useful after the engineering and branch-validity gates have passed.

For the restarted high-dimensional block, the benchmark architecture should at
least cover:
\begin{enumerate}[label=(\arabic*), leftmargin=0.28in]
  \item an exact or controlled linear-Gaussian reference where available,
  \item a nonlinear toy case that reveals posterior-shape or retained-object
  limitations,
  \item a branch-valid finite-difference check on the declared approximate scalar,
  \item structural diagnostics specific to the lane: rank, mass, positivity,
  PSD, support, Jacobian, or ESS as appropriate.
\end{enumerate}
These ladders are not optional appendices.  They are what stops a chapter-local
approximation from being promoted on speed or elegance alone.

The benchmark backbone should be named explicitly so the reader does not have to
reconstruct it from isolated examples:
\begin{enumerate}[label=(V\arabic*), leftmargin=0.34in]
  \item exact-reference linear Gaussian benchmark,
  \item long-horizon stochastic-volatility benchmark,
  \item spatial SIR benchmark for high-dimensional partial-observation filtering,
  \item predator-prey benchmark for nonlinear preconditioning,
  \item derivative validation table whenever a fixed-branch gradient is part of
  the claim.
\end{enumerate}
A useful compact benchmark table should pair the resource ledger with the
model-specific accuracy ledger and, when derivatives are claimed, with a separate
finite-difference table containing \(D_h\), \(G_{\rm FB}\), \(|D_h-G_{\rm FB}|\),
and the branch-status fields for each \(h\).  The derivative table is meaningful
only when the branch record is identical in the three columns
\(\theta-h\), \(\theta\), and \(\theta+h\).

\paragraph{What each benchmark can and cannot establish.}
The linear-Gaussian benchmark can establish exact-reference agreement for the
declared run contract, but it cannot establish nonlinear posterior accuracy.  The
stochastic-volatility benchmark can establish long-horizon nonlinear stability and
credible agreement with references or synthetic truth under the declared run
contract, but it cannot by itself establish high-dimensional scalability.  The
spatial SIR benchmark can establish whether the method remains useful at the
tested latent dimension, horizon, rank ladder, and basis ladder, but it cannot
prove that all spatial graphs or epidemic models admit moderate TT rank.  The
predator-prey benchmark can establish that nonlinear preconditioning yields a
measurable benefit on the tested nonlinear dynamical system, but it cannot prove
that such preconditioning is always worth its cost.  One-seed demonstrations are
diagnostic unless the result artifact explicitly justifies stronger
interpretation.

\section{Defect Calculus And Promotion Rules}
\label{sec:bf-hd-promotion-logic}

The chapter closes with a simple rule: runtime, memory, ESS per gradient, rank,
and point count can compare methods only after the target, semantic, and
numerical veto diagnostics have passed.  The industrial promotion logic is
therefore conditional:
\begin{equation}
  V \text{ passes } \Longrightarrow \text{compare performance};
  \qquad
  V^c \Longrightarrow \text{veto, not ranking.}
  \label{eq:bf-hd-validation-veto-logic}
\end{equation}
This chapter does not conclude that high-dimensional nonlinear filtering is
solved for BayesFilter, that any method is NAWM-ready, that HMC converges, or
that any lane should become a production default.  It concludes only that a
candidate method must pass the declared validation and defect calculus before
speed, elegance, or novelty are allowed to matter.
